\section{Introduction}\label{sec:introduction}

Normalized volume, introduced by Chi Li \cite{Li18}, measures the
size of a klt singularity by balancing
the discrepancy of a valuation against the asymptotic colength of its
valuation ideals. For a closed point \(x\) of an \(n\)-dimensional
normal complex variety that is klt near \(x\), with zero boundary,
it is
\[
 \nvol(x,X)=\inf_{v\text{ centered at }x} A_X(v)^n\vol(v).
\]
Here \(A_X(v)\) is log discrepancy, and \(\vol(v)\) is the
leading colength of the ideals \(\{f:v(f)\ge t\}\), multiplied by
\(n!\). Section~\ref{sec:cones} gives the precise conventions for
real valuations, including infinite discrepancy. In dimension \(n\),
the largest normalized volume is \(n^n\),
attained exactly at smooth points \cite[Theorem~1.6]{LX19}.
The ordinary double point gap conjecture predicts the next largest
value: every singular klt germ should have normalized volume at most
\(2(n-1)^n\), and this value should characterize the ordinary double
point. Besides distinguishing a basic singularity by a numerical
invariant, such a gap gives global restrictions on K-semistable Fano
varieties through the local-to-global volume inequality
\cite[Theorem~21]{Liu18}.

An ordinary double point is the analytic hypersurface germ of a
nondegenerate quadratic form. It is the simplest singular quadratic
cone, and the asserted gap concerns all singular klt germs, not only
hypersurfaces. Thus an essential part of the problem is to extract
structural information from a numerical lower bound on volume.

\Needspace{17\baselineskip}
\begin{theorem}[The ordinary double point gap]\label{thm:main}
Let \(n\ge2\), and let \(x\) be a
singular closed point of a normal complex algebraic \(n\)-fold \(X\).
Suppose that \(K_X\) is \(\QQ\)-Cartier and that \(X\) is klt near \(x\),
with zero boundary. Then
\[
                    \nvol(x,X)\le2(n-1)^n.
\]
Equality holds if and only if
\[
 (X,x)_{\mathrm{an}}\cong
 \left(\{z_0^2+\cdots+z_n^2=0\}\subset\CC^{n+1},0\right).
\]
\end{theorem}

The point need not be isolated in the singular locus. No structural
assumption on the singularity is implicit in Theorem~\ref{thm:main}.
We use log discrepancy, so the discrepancy of the exceptional divisor
of the blowup of a smooth \(n\)-fold point is \(n\). The normalized
volume at the ordinary double point is \(2(n-1)^n\), including its
interpretation as an infimum over all centered valuations.

The proof uses the fourfold theorem of the companion paper
\emph{The normalized-volume gap in dimension four}
\cite[Theorem~1.1]{FourfoldCompanion} as its sole companion input.
That theorem has the same unrestricted boundary-zero scope and includes
the analytic equality characterization; we restate it in
Section~\ref{sec:cones}. The induction from this base case, developed
below, proves both the bound and the equality assertion in every higher
dimension.

\begin{corollary}[Global volume bounds]\label{cor:global}
Let \(V\) be an \(n\)-dimensional
K-semistable \(\QQ\)-Fano variety, where \(n\ge2\).
If \(V\) is singular, then
\[
       (-K_V)^n\le
       2\left(\frac{n^2-1}{n}\right)^n<2n^n.
\]
If \(V\not\cong\PP^n\), then \((-K_V)^n\le2n^n\), with equality
precisely when \(V\) is a smooth quadric \(Q^n\subset\PP^{n+1}\) or
\(\PP^1\times\PP^{n-1}\). These two descriptions coincide for \(n=2\).
\end{corollary}

Here a \(\QQ\)-Fano variety is normal, projective and klt, with ample
\(\QQ\)-Cartier anticanonical divisor. The smaller singular bound in
Corollary~\ref{cor:global} is not accompanied by an assertion of
sharpness or a classification of its equality cases.

\subsection{Earlier results}

The gap conjecture arose from the study of singular noncollapsed
polarized K\"ahler--Einstein limits, where the volume density measures the
size of a metric tangent cone. The algebraicity and uniqueness of
these tangent cones were established by Donaldson--Sun
\cite{DonaldsonSun17}. Spotti--Sun proposed that the ordinary double
point has the largest density among singular Ricci-flat K\"ahler
cones with isolated vertex, and formulated the corresponding
normalized-volume conjecture for arbitrary singular klt germs
\cite[Conjectures~1.2 and~5.5]{SpottiSun17}. Their observation that
density above \(1/2\) forces Gorenstein canonical structure in the
metric-limit setting \cite[Proposition~3.7]{SpottiSun17} already
illustrates a central feature of the problem: a sufficiently large
volume can force algebraic structure that is absent from the
hypotheses. Here that principle is applied after stable degeneration,
starting from an arbitrary klt germ.

The unrestricted low-dimensional cases precede the present induction.
For surfaces, Li--Liu's calculation
\(\nvol(0,\CC^2/G)=4/|G|\) for finite groups acting freely in
codimension one, together with the quotient classification, gives
the bound and its equality case
\cite[Proposition~4.10]{LL19}. Liu--Xu proved the threefold theorem
for arbitrary klt germs \cite[Theorem~1.3]{LX19}.
In higher dimensions, Liu proved the conjecture for local complete
intersections \cite[Theorem~1.3]{Liu22}. For singular \(\QQ\)-Gorenstein toric germs,
Moraga--S\"u\ss\ proved the different sharp bound
\(\nvol\le(16/27)n^n\) for \(n\ge3\), with equality precisely
for the product of a three-dimensional ordinary double point and
\(\CC^{n-3}\); in dimension two the bound is \(2\), with equality
at the surface ordinary double point
\cite[Theorem~2]{MoragaSuess24}. The fourfold companion supplies
the remaining base case used here; the higher-dimensional argument
does not assume either a complete-intersection or a toric structure.

The global application combines the local gap with a complementary
smooth theorem. Li--Miao proved the bound \(2n^n\) and its equality
classification for K-semistable Fano manifolds other than projective
space \cite[Theorem~1.1]{LM25}. Li--Miao--Zhang established the
corresponding metric gap for compact K\"ahler manifolds other than
projective space with \(\operatorname{Ric}(\omega)\ge(n+1)\omega\),
and a global degree bound for K-semistable toric log Fano pairs
whose underlying variety is not projective space
\cite[Theorems~1.3 and~1.6]{LiMiaoZhang26}.
The local theorem here gives a strict improvement for singular
K-semistable \(\QQ\)-Fano varieties, so that equality at \(2n^n\)
reduces to Li--Miao's smooth classification.
Further local refinements address the range below the threefold
maximum: Liu proves the bound \(9\) for Gorenstein canonical
non-hypersurface threefolds and classifies all klt threefold germs
with normalized volume at least \(9\)
\cite[Theorems~1.1 and~1.2]{LiuThreefold26}.

The passage from metric cones to arbitrary klt singularities rests
on the variational theory of normalized volume. Martelli--Sparks--Yau
proved volume minimization among Reeb fields with fixed canonical
weight \cite{MSY08}; Li--Liu and Li--Xu related cone metrics and
K-stability to minimization over all centered valuations
\cite[Theorem~1.9]{LL19}\cite[Theorem~1.3]{LiXu18}.
For general klt germs, existence of a minimizer, its
quasi-monomiality, uniqueness up to scaling, and finite generation
were established by Blum, Xu, and Xu--Zhuang
\cite{Blum18,Xu20,XZ21,XZStable}.
These results produce a K-semistable cone whose Reeb valuation
computes the original normalized volume. The further degeneration
of Li--Wang--Xu gives a K-polystable cone with the same value
\cite[Theorem~1.2]{LWX21}.
Thus the cone reduction preserves the infimum over all valuations,
which is essential for a gap theorem about arbitrary germs.

We also use lower semicontinuity \cite{BL21} and the finite degree
formula \cite{XZ21} to derive smoothness of the puncture and
Gorenstein structure from the density threshold.
The new argument begins after that reduction: it must control
integral gradings approximating the minimizing Reeb vector of an
isolated Gorenstein cone, whose quotient is
usually an orbifold. Passing to the coarse quotient too early loses
the cyclic characters that control both rational curves and the
eventual volume estimate.

Several classical methods enter at this point, in different roles.
The use of rational curves belongs to the bend-and-break approach
initiated by Mori \cite{Mori79}, with stack-theoretic extensions such
as Chen--Tseng \cite{ChenTseng09}. The particular curve needed here
has at most one stack point and comes from Li--Zhou \cite{LZ25}.
Properness of twisted stable maps \cite{AOV11} turns its minimal
degree into a finite evaluation map. The Ricci-flat cone metric
supplied by Collins--Sz\'ekelyhidi \cite{CS19} then bounds the weights
of symmetric tensors obtained from this map. This is related in
method to Reeb-weight obstructions for holomorphic functions
\cite{GMSY07}, but the symmetric-tensor estimate is proved here.
The final high-index classification is the projective-space and
quadric characterization of Kobayashi--Ochiai \cite{KO73}; the
needed ample-line-bundle formulation also follows from
Araujo--Druel--Kov\'acs \cite{ADK08}.

The other case uses an adjoint construction. Howald's multiplier-ideal
formula \cite{Howald01} and Kawamata--Viehweg vanishing
\cite{FujinoKV} lift a prescribed leading Taylor monomial after the
canonical shift. Hochster's theorem and Stanley's canonical-module
formula \cite{Hochster72,Stanley78} subsequently convert the resulting
normal semigroup into a Gorenstein degeneration. None of those
statements alone forces the canonical point to be interior: that
separate numerical assertion is the content of the exponential
estimate in Section~\ref{jet:section}. Its semigroup counting follows
the convex-geometric approach to graded linear series developed by
Khovanskii, Okounkov, Lazarsfeld--Musta\c t\u a, and
Kaveh--Khovanskii
\cite{Khovanskii1992,Okounkov2003,LazarsfeldMustata2009,KavehKhovanskii2012}.
The new exact saturation conclusion is stronger than the asymptotic
count supplied by that method and requires the adjoint argument.

\subsection{The argument and its additional statements}

Put \(n=N+1\) and
\[
 p_n=2\left(\frac{n-1}{n}\right)^n,\qquad
 d(x,X)=\frac{\nvol(x,X)}{n^n}.
\]
Starting from the known cases in dimensions two and three and the
fourfold companion, we argue by induction. In the step \(n\ge5\), for a
singular germ with
\(d(x,X)\ge p_n\), stable degeneration preserves this density and
bounds the original embedding dimension by that of a K-polystable
cone \(C\). The lower-dimensional gap and the finite degree formula
force \(C\) to be smooth off its vertex and Gorenstein. These are
conclusions of the reduction, rather than assumptions on the germ.

A Reeb vector is a real torus grading positive on every nonconstant
homogeneous function. Approximate the minimizing Reeb ray of \(C\)
by primitive integral gradings. For one such grading, let \(r\) be
the weight of a generating canonical form and let \(m_{\max}\) be
the largest stabilizer order on the puncture. The quotient by the
grading action is a smooth orbifold \(B\), with faithful ample
weight-one line bundle \(L\) and \(-K_B=rL\).
The comparison of \(m_{\max}\) with \(r/N\) separates the two
parts of the argument.

When \(m_{\max}\le r/N\), a minimal orbifold rational curve
constrains the tangent characters. Unless those constraints already
force the grading to be free, its deformations give a finite evaluation
map onto \(C\). Taking the product of their infinitesimal translations
over the generic sheets produces a nonzero symmetric tensor. A weight
bound from the Ricci-flat cone metric then forces a free grading with
canonical weight \(N\) or \(N+1\). The high-index classification
identifies \((B,L)\) as a quadric or projective space with its standard
polarization, so \(C\) is an ODP or is smooth.

For \(m_{\max}\ge r/N\), choose a point with maximal stabilizer
and record each section's degree and leading transverse Taylor exponent.
These pairs form a semigroup \(\Gamma\) in the lattice determined
by the stabilizer characters; write \(K\) for its closed convex cone.
The canonical form determines the lattice point
\(c_0=(r,1,\ldots,1)\). Two independent arguments give exact control
of \(\Gamma\). Near-minimality of the Reeb grading gives moment
inequalities, and an exponential expectation bound puts \(c_0\) in
\(\Int K\). Adjoint vanishing lifts Taylor monomials whose indices,
after the shift by \(c_0\), lie in that interior. Together these facts
force rational polyhedrality and exact saturation of \(\Gamma\).

Two filtrations convert this semigroup structure into a singular klt
slice \(Y\) of the same dimension \(n\). A nontrivial stabilizer
of order \(h\ge2\) for the second filtration action makes one
character account for only \(1/h\) of the leading valuation
colength on \(Y\). This improves the volume comparison with \(C\).
The action and its order \(h\) are different from those defining
\(m_{\max}\); the quadratic-cone example in
Section~\ref{filt:section} makes the distinction explicit.

To use this same-dimensional slice without assuming the desired bound
for it, let \(M_n\) be the supremum of densities of singular
boundary-zero klt \(n\)-fold germs. The comparison gives
\[
 d(o,C)\le\frac{p_n+M_n}{2}.
\]
Applying this estimate to densities approaching \(M_n\) proves
\(M_n\le p_n\). The approximation of gradings is made separately
on each fixed cone; the supremum need not be attained. In the equality
case, a strict colength count and integral rounding return the cone
to the small-stabilizer case. Its ODP embedding dimension then makes
the original singular germ a hypersurface, and a weighted test forces
a nondegenerate quadratic part. Figure~\ref{fig:route} records this
case split and the return in equality.

\begin{figure}[tb]
\centering
\begin{tikzpicture}[
  proofbox/.style={draw=black!65,rounded corners=2pt,align=center,
    inner sep=6pt,font=\small,text width=5.4cm},
  proofarrow/.style={-{Stealth[length=2mm]},draw=black!75},
  node distance=8mm and 8mm]
\node[proofbox,text width=11.6cm] (cone)
  {\(n\ge5\), \(d(x,X)\ge p_n\): stable degeneration and induction\\
   An isolated Gorenstein K-polystable cone \(C\)};
\node[proofbox,below=9mm of cone.south,xshift=-31mm] (small)
  {\(m_{\max}\le r/N\)\\
   Minimal curve: direct rigidity\\
   or finite evaluation and tensor estimate};
\node[proofbox,below=9mm of cone.south,xshift=31mm] (large)
  {\(m_{\max}\ge r/N\)\\
   Canonical point \(c_0\in\Int K\)\\
   Adjoint lifting and saturation};
\node[proofbox,below=8mm of small] (classify)
  {Smooth projective quotient\\
   \((\PP^N,\cO(1))\) or \((Q^N,\cO(1))\)\\
   Cone smooth or ODP};
\node[proofbox,below=8mm of large] (slice)
  {Two filtrations and a singular slice \(Y\)\\
   Character share \(1/h\), \(h\ge2\)};
\node[proofbox,below=8mm of slice] (bound)
  {\(d(o,C)\le (p_n+M_n)/2\)\\
   Hence \(M_n\le p_n\)};
\draw[proofarrow] (cone.south) -- ++(0,-4mm) -| (small.north);
\draw[proofarrow] (cone.south) -- ++(0,-4mm) -| (large.north);
\draw[proofarrow] (small) -- (classify);
\draw[proofarrow] (large) -- (slice);
\draw[proofarrow] (slice) -- (bound);
\coordinate (returnleft) at ($(small.west)+(-7mm,0)$);
\coordinate (returnbase) at ($(bound.south)+(0,-7mm)$);
\coordinate (returnturn) at (returnleft |- returnbase);
\draw[proofarrow,dashed] (bound.south) -- (returnbase)
  -- node[midway,below,font=\scriptsize]
     {Equality and integral rounding} (returnturn)
  -- (returnleft) -- (small.west);
\end{tikzpicture}

\caption{The induction step \(n=N+1\ge5\) for a singular germ of
density at least \(p_n\). Here \(K\) is the cone of leading Taylor
indices and \(c_0\) its canonical lattice point.
The same-dimensional slice enters only through the supremum \(M_n\)
of singular \(n\)-dimensional densities.
Equality in the volume estimate returns the cone to the small
stabilizer case.}
\label{fig:route}
\end{figure}

\paragraph{Organization.}
Section~\ref{sec:cones} fixes the valuation conventions, supplies the
cone reduction, and proves the absolute grading approximation.
Section~\ref{small:section} treats small stabilizers. For the other
case, Section~\ref{jet:section} proves canonical interiority and
Section~\ref{adj:section} independently establishes adjoint lifting
before combining the two results. Section~\ref{filt:section}
constructs the two filtrations and the second stabilizer.
Section~\ref{boot:section} constructs the same-dimensional singular
slice, proves the supremum inequality, and analyzes equality.
Finally, Section~\ref{sec:conclusion} returns to the original germ
and derives the Fano consequences.
